% TOP_PROBLEM: {"id": 1245, "title": "Newman's Conjecture", "category_id": 1, "category": {"id": 1, "name": "number_theory", "display_name": "Number Theory", "description": "Properties of integers, prime numbers, Diophantine equations.", "slug": "number-theory", "order_index": 1, "created_at": "2026-07-31T15:26:25.670Z"}, "status": "open"}
% FIRST_POSTED: null
% ENTRY_KIND: statement_only
% Imported from: batch05.tex; UnsolvedMath ID 1245
\documentclass[11pt]{article}
\usepackage[margin=25mm]{geometry}
\usepackage{fontspec}
\setmainfont{Latin Modern Roman}
\usepackage{amsmath,amssymb,mathtools}
\usepackage{unicode-math}
\setmathfont{Latin Modern Math}
\usepackage{xurl,hyperref}
\hypersetup{colorlinks=true,urlcolor=blue}
\setcounter{secnumdepth}{0}
\setlength{\parindent}{0pt}
\setlength{\parskip}{5pt}
\setlength{\emergencystretch}{3em}
\newcommand{\sref}[2]{\hyperlink{source:#1:#2}{[S#2]}}
\newcommand{\cataloguescope}[1]{\paragraph{Recorded target.}#1}
\begin{document}
% UnsolvedMath ID 1245; source catalogue code NT-038
\setcounter{section}{1244}
\section{Newman's Conjecture}\label{1245}
{\small\sffamily UnsolvedMath ID 1245 · Number Theory}
\subsection{Definitions and mathematical statement}

\paragraph{Shared notation.}$\mathbb N=\{1,2,\ldots\}$, and $\mathbb Z,\mathbb Q,\mathbb R,\mathbb C$ denote the integers, rationals, reals and complex numbers. Any other conventions are specified locally.


For $n\geq1$, let $p(n)$ be the number of integer partitions of $n$, that is, the number of finite nonincreasing sequences of positive integers with sum $n$. Set $p(0)=1$. For integers $m\geq1$ and $a$, a unit residue means $\gcd(a,m)=1$.
\paragraph{Statement recorded in the supplied source.}
For every integer $m\geq1$ and every integer $a$ with $\gcd(a,m)=1$, is the set
\[
\{n\in\mathbb N:p(n)\equiv a\pmod m\}
\]
infinite? \sref{1245}{2}
The usual Newman conjecture asks the same question for every residue $a$, including nonunits. The recorded unit-residue question is the weaker restriction appearing in this catalogue\textquotesingle s background. \sref{1245}{2}
\subsection{Short English statement}
Does every residue coprime to its modulus occur infinitely often among partition numbers?
\subsection{Sources}
\begin{description}
\item[\hypertarget{source:1245:1}{S1}] ULAM.AI and the UnsolvedMath Contributors, UnsolvedMath: A Curated Collection of Open Mathematics Problems, record 1245 (NT-038). CC BY 4.0. \url{https://huggingface.co/datasets/ulamai/UnsolvedMath}.
\item[\hypertarget{source:1245:2}{S2}] Dohoon Choi and Youngmin Lee, \emph{Newman\textquotesingle s conjecture for the partition function modulo integers with at least two distinct prime divisors}, arXiv:2212.00636v2 (2025), Section 1, Conjecture 1.1. \url{https://arxiv.org/pdf/2212.00636}.
\end{description}
\label{end:1245}
\end{document}
